\PassOptionsToPackage{table,xcdraw}{xcolor}
\documentclass[11pt, a4paper]{article}

\usepackage[T1]{fontenc}
\usepackage[utf8]{inputenc}
\usepackage{tgpagella}
\usepackage[spanish, es-noshorthands]{babel}
\usepackage{amsmath, amssymb}
\usepackage[most]{tcolorbox}
\usepackage{tabularx}
\usepackage[margin=2.5cm]{geometry}
\usepackage{fancyhdr}

\pagestyle{fancy}
\fancyhf{}
\setlength{\headheight}{16pt}
\lhead{\textbf{UCB Sede Tarija} | Ing. Mecatrónica}
\rhead{IMT-221 | Rúbrica Analítica de Calificación}
\renewcommand{\headrulewidth}{0.4pt}
\definecolor{ucbblue}{RGB}{0, 51, 102}

\begin{document}

\begin{tcolorbox}[colback=ucbblue!5, colframe=ucbblue, title=\textbf{Rúbrica Analítica de Evaluación — EC2 (BJT / Par Diferencial)}]
    Evalúe el razonamiento del estudiante en base a 5 dimensiones transversales, sin basar la puntuación únicamente en el resultado aritmético final[cite: 8].
\end{tcolorbox}

\renewcommand{\arraystretch}{1.8}
\begin{tabularx}{\textwidth}{|>{\raggedright\arraybackslash}p{2.5cm}|>{\raggedright\arraybackslash}X|>{\raggedright\arraybackslash}X|>{\raggedright\arraybackslash}X|}
\hline
\rowcolor{ucbblue!20}
\textbf{Dimensión} & \textbf{Strong (Sólido)[cite: 8]} & \textbf{Developing (En desarrollo)[cite: 8]} & \textbf{Weak (Débil)[cite: 8]} \\ \hline
\textbf{1. Selección del Modelo} & Elige el modelo correcto (ej. LVK de colector, Espejo ideal/$\beta$ finito) justificando sus supuestos teóricos explícitamente. & Usa el modelo mayormente correcto pero con justificaciones débiles o incompletas de las condiciones frontera. & Uso mecánico de fórmulas desvinculadas de la topología real del circuito evaluado. \\ \hline
\textbf{2. Ejecución Analítica} & Las ecuaciones, despejes, el manejo de unidades y la aritmética final son coherentes y libres de errores. & Demuestra el camino analítico correcto pero comete errores menores (de signo o conversión de mV a V). & El camino central de la derivación analítica o la relación matemática básica es incorrecto. \\ \hline
\textbf{3. Interpretación Física} & Logra explicar el significado físico de sus resultados numéricos (ej. por qué el $\beta$ no explica toda la caída de $I_T$). & Proporciona interpretaciones parciales o repite definiciones genéricas que no calzan exacto con sus números. & Proveé una respuesta numérica final en vacío, sin capacidad de conectar el número a la física del circuito. \\ \hline
\textbf{4. Consistencia de Convención} & Mantiene rigurosamente la convención exigida (Ej. $v_o=v_{C2}$) a lo largo de los cálculos de $A_d, A_c$ y CMRR. & Mezcla temporalmente nomenclaturas pero nota la inconsistencia tras una leve revisión procedimental. & Mezcla convenciones (Ej. Divide \textit{diff} entre \textit{single-ended}) sin ser consciente del error metodológico. \\ \hline
\textbf{5. Diagnosis \& Troubleshooting} & Propone mediciones específicas que prueban directa y lógicamente una hipótesis de falla de hardware plausible. & La medición propuesta es relevante pero está débilmente conectada con la causa raíz hipotética. & Propone respuestas arbitrarias, adivinanzas o "cambiar componentes al azar" sin usar instrumentación. \\ \hline
\end{tabularx}

\vspace{0.5cm}
\textit{Nota al Docente: Según las directrices actuales, no se imponen ponderaciones numéricas formales en este artefacto; debe ajustarlas acorde al sistema de la universidad[cite: 8].}

\end{document}
